\section{Design}

\subsection{Design Goals}
The design of the system follows one central principle: each major responsibility
should be owned by one RTOS task, and communication between tasks should happen
through explicit synchronization primitives rather than implicit shared-state
dependencies. This is especially important for this project because the system
combines periodic sensing, asynchronous user interaction, local rendering,
network-state transitions, TinyML inference, and MQTT communication. If these
concerns were implemented in a single control loop, the result would be harder
to reason about and harder to evaluate as an RTOS design.

The architecture therefore aims to satisfy five goals:
\begin{enumerate}
  \item predictable periodic sensing;
  \item clear separation between producers and consumers of data;
  \item minimal blocking between unrelated modules;
  \item controlled access to shared resources such as the I\textsuperscript{2}C
  bus and shared state; and
  \item graceful transition between AP mode, STA mode, local UI, and cloud
  connectivity.
\end{enumerate}

\subsection{System Architecture Overview}
At the highest level, the system can be viewed as a streaming pipeline centered
around sensor acquisition. The \texttt{sensor\_task} periodically produces the
latest environmental measurements and distributes them to the rest of the
system. Several downstream tasks consume the same information for different
purposes:
\begin{itemize}
  \item the LED task converts temperature into blink behavior,
  \item the NeoPixel task converts humidity into color,
  \item the LCD task renders local status,
  \item the web task updates the browser-visible dashboard and history buffer,
  \item the TinyML task performs prediction, and
  \item the CoreIOT task publishes telemetry to the cloud.
\end{itemize}

The web task is also the owner of Wi-Fi mode management. It decides whether the
board is running in AP mode or attempting STA connection, and it signals the
CoreIOT task when Internet connectivity becomes available. The TinyML and
CoreIOT tasks are therefore not directly responsible for network setup; they
are consumers of network availability.

\subsection{Hardware Mapping}
Table~\ref{tab:hardware-mapping} summarizes the active hardware resources used
by the project.

\begin{table}[H]
\centering
\caption{Hardware mapping of the implemented system}
\label{tab:hardware-mapping}
\begin{tabular}{p{3.2cm}p{2.5cm}p{8.0cm}}
\toprule
\textbf{Component} & \textbf{Pin / Address} & \textbf{Purpose} \\
\midrule
ESP32-S3 board & -- & Main microcontroller, Wi-Fi connectivity, RTOS execution \\
DHT20 sensor & I\textsuperscript{2}C bus & Temperature and humidity acquisition \\
LCD 16x2 & I\textsuperscript{2}C address \texttt{0x21} & Local system display for sensor, Wi-Fi, and TinyML states \\
I\textsuperscript{2}C SDA & GPIO11 & Shared bus data line \\
I\textsuperscript{2}C SCL & GPIO12 & Shared bus clock line \\
Discrete LED & GPIO48 & Temperature-dependent blink feedback \\
NeoPixel & GPIO45 & Humidity-dependent color indication \\
External LED & GPIO5 & User-controlled LED from the web interface, PWM brightness \\
BOOT button & GPIO0 & Manual fallback from STA/AP flow back to AP mode \\
\bottomrule
\end{tabular}
\end{table}

The mapping shows that the architecture is constrained by both physical sharing
and functional reuse. The DHT20 and LCD share the I\textsuperscript{2}C bus,
which justifies the use of a mutex. The external LED is distinct from the
temperature LED, which makes it possible to satisfy the web control task
without interfering with the environmental blink behavior.

\subsection{RTOS Task Architecture}
The complete runtime is built from seven explicitly created tasks. Their
responsibilities, priorities, and triggering models are shown in
Table~\ref{tab:task-architecture}.

\begin{table}[H]
\centering
\caption{FreeRTOS task architecture}
\label{tab:task-architecture}
\begin{tabular}{p{3cm}p{1cm}p{1.8cm}p{3.5cm}p{5.2cm}}
\toprule
\textbf{Task} & \textbf{Prio.} & \textbf{Stack} & \textbf{Activation model} & \textbf{Main responsibility} \\
\midrule
\texttt{sensor\_task} & 3 & 4096 & Periodic every 2 s & Read DHT20, validate values, dispatch samples to downstream consumers \\
\texttt{lcd\_task} & 2 & 4096 & Event-driven, queue wait up to 250 ms & Render local LCD content depending on sensor, Wi-Fi, and TinyML state \\
\texttt{main\_server\_task} & 4 & 8192 & Fast service loop every 20 ms & Run HTTP server, handle AP/STA state, Wi-Fi profiles, dashboard logic \\
\texttt{led\_manager\_task} & 1 & 3072 & Event-driven plus timeout & Generate discrete LED blink patterns from temperature or event commands \\
\texttt{neo\_pixel\_task} & 1 & 3072 & Event-driven, blocking queue receive & Map humidity to RGB color and update NeoPixel \\
\texttt{tiny\_ml\_task} & 1 & 8192 & Event-driven by sensor queue & Accumulate samples, build features, invoke TFLM model, publish result \\
\texttt{coreiot\_task} & 1 & 6144 & Event-driven after Internet semaphore + periodic publish & Connect MQTT client, serve RPC, publish telemetry every 10 s \\
\bottomrule
\end{tabular}
\end{table}

This division is a good fit for the project requirements. The web task receives
the highest priority because it is responsible for configuration and network
mode transitions. The sensor task receives the next-highest priority because it
defines the cadence of environmental sampling. Presentation and cloud tasks are
lower priority because they consume already acquired data rather than producing
the primary physical measurements.

\subsection{Application Context and Shared State}
The project uses an \texttt{app\_context\_t} structure as the central ownership
object for all queues, mutexes, the Internet-ready semaphore, and a compact
shared state block. This is an important architectural decision because it
reduces the need for project-wide global variables. Instead of exposing
individual globals for temperature, humidity, Wi-Fi connection state, or LCD
content mode, the code stores them inside a shared state structure protected by
\texttt{stateMutex}.

The shared state currently contains:
\begin{itemize}
  \item latest temperature,
  \item latest humidity,
  \item a boolean Wi-Fi-connected flag, and
  \item the current LCD content mode.
\end{itemize}

This does not remove every file-local static variable from the repository. The
web server module, for example, still maintains private runtime variables for
its internal HTML state, history buffer, and Wi-Fi workflow. However, it
significantly reduces uncontrolled cross-module coupling and improves the
clarity of synchronization.

\subsection{Inter-Task Communication Model}
The project uses three primary synchronization mechanisms.

\subsubsection*{Queues.}
Queues are used for directional data flow. The design intentionally favors
single-slot queues for continuously refreshed information. This means the system
does not attempt to preserve every historical sensor sample inside RTOS queues;
instead, each consumer receives the most recent snapshot. This choice is
appropriate for UI and cloud applications where stale data are less useful than
fresh data.

\subsubsection*{Mutexes.}
Two mutexes are central to the design:
\begin{itemize}
  \item \texttt{i2cMutex} protects the shared I\textsuperscript{2}C bus between
  DHT20 reads and LCD writes.
  \item \texttt{stateMutex} protects access to the shared state block.
\end{itemize}
Without \texttt{i2cMutex}, concurrent sensor and LCD access could lead to
corrupted transactions or inconsistent device behavior. Without
\texttt{stateMutex}, Wi-Fi or LCD content state could be modified while another
task reads it.

\subsubsection*{Binary semaphore.}
The \texttt{internetSemaphore} is used as a one-shot signal from the web task to
the CoreIOT task. It encodes an event rather than a data value: ``STA
connectivity is now ready for cloud communication.'' This is an elegant use of
a binary semaphore because the CoreIOT task should not try to initialize MQTT
before the device is in the correct network mode.

\subsection{Queue Topology}
The queue topology is intentionally asymmetric. The sensor task is the main
producer and fans out data into multiple destinations. Table
\ref{tab:queue-topology} shows the purpose of each queue.

\begin{table}[H]
\centering
\caption{Queue topology and semantics}
\label{tab:queue-topology}
\begin{tabular}{p{3cm}p{2cm}p{4cm}p{6cm}}
\toprule
\textbf{Queue} & \textbf{Length} & \textbf{Producer} & \textbf{Consumer / Purpose} \\
\midrule
\texttt{ledQueue} & 8 & sensor task or system events & LED manager receives mode-change and pulse commands \\
\texttt{neoQueue} & 1 & sensor task & NeoPixel receives latest humidity snapshot \\
\texttt{lcdQueue} & 1 & sensor task & LCD receives latest environmental sample \\
\texttt{webQueue} & 1 & sensor task & Web server receives latest sample and updates history buffer \\
\texttt{coreQueue} & 1 & sensor task & CoreIOT receives latest sample for periodic telemetry \\
\texttt{tinyMLQueue} & 1 & sensor task & TinyML receives latest sample stream for batching and inference \\
\texttt{tinyResultQueue} & 1 & TinyML task & LCD, web, and CoreIOT consume the latest inference result \\
\bottomrule
\end{tabular}
\end{table}

The use of \texttt{xQueueOverwrite()} on most single-slot queues is particularly
appropriate because the system is monitoring-oriented. The most recent
temperature or humidity value is more important than the complete backlog of old
samples.

\subsection{State Flow Through the System}
Figure generation is not embedded directly inside this report, but the
architecture can be reproduced externally using the following Mermaid code.
This code can be pasted into draw.io through
\textit{Arrange $\rightarrow$ Insert $\rightarrow$ Advanced $\rightarrow$
Mermaid}, or into any Mermaid-compatible viewer.

\begin{verbatim}
flowchart LR
    DHT20[DHT20 Sensor] --> ST[Sensor Task]
    ST -->|temperature| LQ[ledQueue]
    ST -->|humidity| NQ[neoQueue]
    ST -->|latest sample| LCDQ[lcdQueue]
    ST -->|latest sample| WQ[webQueue]
    ST -->|latest sample| CQ[coreQueue]
    ST -->|latest sample| TQ[tinyMLQueue]

    LQ --> LT[LED Manager Task]
    NQ --> NT[NeoPixel Task]
    LCDQ --> LCDT[LCD Task]
    WQ --> WT[Web Server Task]
    TQ --> TT[TinyML Task]
    CQ --> CT[CoreIOT Task]

    TT -->|tinyml_result| TRQ[tinyResultQueue]
    TRQ --> LCDT
    TRQ --> WT
    TRQ --> CT

    WT -->|internetSemaphore| CT
    WT --> WIFI[AP/STA Wi-Fi Manager]
    WT --> PREFS[Preferences Storage]
    WT --> BROWSER[HTTP Dashboard]
    CT --> CLOUD[CoreIOT MQTT]
\end{verbatim}

The diagram highlights the main design pattern of the system: one producer task
feeds multiple consumer tasks, while secondary tasks such as TinyML and Wi-Fi
management enrich the data path with derived information and network events.

\subsection{Design Rationale}
Several implementation choices deserve emphasis:
\begin{itemize}
  \item \textbf{Single responsibility per task.} Each task has one primary
  concern, which keeps the code easier to analyze and debug.
  \item \textbf{Latest-value semantics.} For monitoring and display tasks, the
  newest sample is more important than preserving every historical sample in a
  queue.
  \item \textbf{Bounded synchronization.} Critical resources are protected with
  explicit mutexes instead of ad-hoc timing assumptions.
  \item \textbf{Event-based cloud startup.} The CoreIOT task only starts
  MQTT-related work after the web task signals that STA connectivity is valid.
  \item \textbf{Batch-based inference.} TinyML is intentionally slower than raw
  sensing because it aggregates ten samples before making a prediction. This
  lowers noise and creates a more stable high-level output.
\end{itemize}

Overall, the design is a balanced soft real-time architecture. It remains
straightforward to implement and maintain while still providing a clear RTOS
decomposition and enough modularity for future extension.
